\section{Limitações e Itens Pendentes}

Ficam explicitamente registrados como pendentes, por dependerem de ensaio
físico ainda não realizado e não por limitação de escopo. Todos estão
programados para o segundo semestre de \num{2026}, dentro do prazo de
fechamento do trabalho:

\begin{itemize}
    \item Identificação de $C_T$ e $C_Q$ da hélice em bancada, e sua
    comparação com as referências APC e UIUC
    (\autoref{sec:resultado_coef_aero_bancada});
    \item Repetição da calibração do tubo de Pitot registrando pressão e
    temperatura do ambiente, para fixar a sensibilidade $S$ sem a
    ambiguidade da densidade adotada (\autoref{sec:calibracao_pitot});
    \item Montagem final do conjunto motopropulsor na estrutura da bancada,
    com hélice acoplada e cabeamento definitivo
    (\autoref{sec:resultado_montagem_final});
    \item Validação completa da cadeia de medição da SIBA --- bateria,
    células de carga (\autoref{sec:resultado_validacao_siba});
    \item Piso de ruído dos canais da SAME com o ESC chaveando e
    verificação da cadeia de tensão contra instrumento de referência --- o
    ensaio com o motor girando mediu as tensões de fase direto pelo
    osciloscópio, sem passar pelo divisor da placa
    (\autoref{sec:resultado_calibracao_canal_same});
    \item Calibração do ganho do sensor de corrente na topologia de uso, com
    fonte capaz de impor correntes da ordem da faixa nominal de
    $\pm\SI{60}{\ampere}$ --- a caracterização atual cobre a faixa magnética
    completa por meio de enrolamentos auxiliares, mas a dispersão de
    $\pm\SI{9,6}{\percent}$ entre configurações mostra que o ganho da espira
    única não é dedutível das demais (\autoref{tab:hall_result}).
\end{itemize}

Cada um desses itens tem seção e figura ou tabela já reservadas no texto, de
modo que o fechamento do trabalho consiste em executar os ensaios e preencher
esses espaços, sem alteração estrutural do documento.
